\section{问题定义补充}\label{sec:supp-problem}

正文 \S\ref{sec:formulation} 以精简的定义建立了 P1/P2/P3 的评估框架。本节补充实验复现所需的完整形式化约定，包括 DAG 字段、评价指标的公式细节以及 spill 机制的精确参数。

\subsection{DAG 与缓冲生命周期}\label{sec:supp-input}

正文仅给出节点两分类的概要，此处列出完整的字段定义。输入是一个有向无环图 $\daggraph$。节点分为两类：

\begin{enumerate}
  \item 操作节点携带 $(\ksField{Id},\ksField{Op},\ksField{Pipe},\ksField{Cycles},\ksField{Bufs})$，表示在某条硬件管线上的计算或数据搬运指令；\ksField{Bufs} 给出该指令读写的逻辑缓冲。
  \item 缓存管理节点携带 $(\ksField{Id},\ksField{Op},\ksField{BufId},\ksField{Size},\ksField{Type})$，其中 $\mathrm{Op}\in\{\ksOp{alloc},\ksOp{free}\}$。每个逻辑缓冲 $b$ 有唯一分配节点 $a_b$、唯一释放节点 $f_b$、大小 $s_b$ 与所属 cache 类型 $\tau(b)$。
\end{enumerate}

边表示必须满足的执行依赖。忽略缓存管理节点后，样例图通常从 \ksOp{copy-in} 进入，经片上计算或搬运，最终由 \ksOp{copy-out} 写回片外；但本文算法不依赖具体算子族或 motif 名称。调度 $S=(v_1,\ldots,v_N)$ 是覆盖所有必需节点的一条拓扑序。

\subsection{三个评估视角}\label{sec:supp-objective}

正文将 P1/P2/P3 浓缩为一段文字，此处展开其数学定义。三个子问题在同一图上逐步增加约束。令节点 $v$ 对缓冲驻留量的增量为
\begin{equation}\label{eq:supp-mv}
  M_{v}=
  \begin{cases}
    +\,s_b, & v=a_b\ (\ksOp{alloc}),\\
    -\,s_b, & v=f_b\ (\ksOp{free}),\\
    0, & \text{其他节点}.
  \end{cases}
\end{equation}

\paragraph{\mdseries P1：峰值驻留。}
P1 只评判节点序，不要求物理地址或 spill 列表。以前缀和定义
\begin{equation}\label{eq:supp-vstay}
  V_{\mathrm{stay}}(0)=0,\qquad
  V_{\mathrm{stay}}(i)=V_{\mathrm{stay}}(i-1)+M_{v_i},\qquad
  V_{\mathrm{stay}}^{\max}=\max_{0\le i\le N}V_{\mathrm{stay}}(i).
\end{equation}
主目标为最小化 $\max V_{\mathrm{stay}}$。本文在报告中同时给出逐 cache 峰值，便于解释 P1 序与 P2/P3 之间的错配。

\paragraph{\mdseries P2：额外 DDR 搬运量。}
P2 在容量约束下给出物理地址和 spill 列表，目标为最小化 spill 引入的额外数据搬运
\begin{equation}\label{eq:supp-extra}
  \extra(S)=\sum_{b\in \mathrm{Spilled}(S)} e_b,\qquad
  e_b=\kappa_b s_b,\qquad
  \kappa_b=
  \begin{cases}
  1, & b\ \text{为 clean},\\
  2, & b\ \text{为 dirty}.
  \end{cases}
\end{equation}

\paragraph{\mdseries P3：流水执行时间。}
P3 在插入 spill 节点和地址复用依赖后，按多管线最早发射模型计算 makespan。记起始时间 $\sigma(v)$、结束时间 $\epsilon(v)$、增强依赖集合为 $E'$：
\begin{equation}\label{eq:supp-time}
  \sigma(v)\ge 0,\quad
  \sigma(v)\ge \epsilon(u)\ \ \forall (u,v)\in E',\quad
  \epsilon(v)=\sigma(v)+\mathrm{Cycles}(v),\quad
  T=\max_v \epsilon(v).
\end{equation}
缓存管理节点周期为 $0$；同一 \ksField{Pipe} 上的节点按调度序串行执行，不允许在管线内部重新排序。

\paragraph{词典键。}
本文以与正文一致的词典键进行比较：P1 使用 $\bigl(\max\livec_{\mathrm{L1}},\max\livec_{\mathrm{UB}},\mathrm{L0\ count},T\bigr)$；P2 使用 $\bigl(\extra,\#\mathrm{spills},T\bigr)$；P3 使用 $\bigl(T,\extra,\#\mathrm{spills}\bigr)$。基准定义未给出将运行时间与额外搬运合并的单一标量，因此 P3 的次级目标在本文中显式采用该词典序口径。

\subsection{容量、地址与 spill 规则}\label{sec:supp-constraints}

正文简述了 clean/dirty 的非对称代价；此处给出地址分配与 spill 机制的完整规则。P2/P3 中，缓冲 $b$ 以偏移 $o_b$ 占据同一 cache 内的连续区间 $[o_b,o_b+s_b)$。任一时刻同时驻留且属于同一 cache 的缓冲区间不得相交，并且必须满足 $0\le o_b$、$o_b+s_b\le \Capc_{\tau(b)}$。容量参数见表~\ref{tab:supp-cap}。

\begin{table}[htbp]
\centering
\caption{硬件 cache 类型与容量（抽象单位）。}\label{tab:supp-cap}
\begin{tabular}{lccccc}
\toprule
Cache 类型 & L1 & UB & L0A & L0B & L0C \\
\midrule
容量 $\Capc$ & 4096 & 1024 & 256 & 256 & 512 \\
\bottomrule
\end{tabular}
\end{table}

本文称缓冲 $b$ 为 clean，当且仅当它出现在某个 \ksOp{copy-in} 节点的 \ksField{Bufs} 中；此时片外已有备份，换出无需写回。其余缓冲称为 dirty。相应的 spill 周期模型为
\begin{equation}\label{eq:supp-spillcycle}
  \mathrm{Cycles}(\mathrm{SPILL\_IN}_b)=2s_b+150,\qquad
  \mathrm{Cycles}(\mathrm{SPILL\_OUT}_b)=
  \begin{cases}
    0, & b\ \text{为 clean},\\
    2s_b+150, & b\ \text{为 dirty}.
  \end{cases}
\end{equation}

若原图有 $N$ 个节点，第 $k$ 次 spill 生成一对隐式节点：
\begin{equation}\label{eq:supp-spillid}
  \mathrm{Id}(\mathrm{SPILL\_OUT}_k)=N+2(k-1),\qquad
  \mathrm{Id}(\mathrm{SPILL\_IN}_k)=N+2(k-1)+1.
\end{equation}
\ksOp{spill-out} 使用 MTE3，\ksOp{spill-in} 使用 MTE2。目标缓冲 $b$ 的基本依赖链为
\begin{equation}
  a_b\to \mathrm{SPILL\_OUT}\to \mathrm{SPILL\_IN}\to f_b.
\end{equation}
已经使用过 $b$ 的操作节点必须位于 \ksOp{spill-out} 之前，尚未使用 $b$ 的操作节点必须位于 \ksOp{spill-in} 之后。物理地址可在 \ksOp{spill-in} 时改变，因此可行性按驻留段检查，而不是只按逻辑生命期 $a_b\to f_b$ 检查。

地址复用会引入额外依赖：若缓冲 $b$ 复用更早缓冲 $a$ 的物理区间，则计时前加入 $f_a\to a_b$。该依赖不改变 spill 流量，却可能改变 P3 makespan。

\subsection{Benchmark 数据描述}\label{sec:supp-data}

本文使用的六个公开 NPU 核内调度实例规模见表~\ref{tab:supp-cases}。所有实验均从 DAG 图结构出发，不使用针对 Conv、FlashAttention 或 Matmul 的专用规则。

\begin{table}[htbp]
\centering
\caption{公开 NPU 核内调度实例规模。}\label{tab:supp-cases}
\begin{tabular}{lrrrr}
\toprule
任务 & 节点数 & 边数 & 缓冲数 & 操作节点数 \\
\midrule
Matmul\_Case0         &  4160 &  7104 &  1216 &  1728 \\
Matmul\_Case1         & 30976 & 55040 &  8960 & 13056 \\
FlashAttention\_Case0 &  1716 &  2712 &   572 &   572 \\
FlashAttention\_Case1 &  6952 & 11184 &  2328 &  2296 \\
Conv\_Case0           &  2580 &  3869 &   831 &   918 \\
Conv\_Case1           & 36086 & 85653 & 12013 & 12060 \\
\bottomrule
\end{tabular}
\end{table}
